\exercisesheader{}

% 31

\eoce{\qt{Pelanggan kedai kopi\label{coffee_shop_customers}} Sebuah kedai kopi
melayani rata-rata 75 pelanggan per jam pada jam sibuk pagi hari.
\begin{parts}
\item
  Distribusi manakah yang telah kita pelajari yang paling tepat
  untuk menghitung peluang datangnya sejumlah tertentu pelanggan
  dalam waktu satu jam
  pada waktu seperti ini?
\item Berapakah rata-rata dan simpangan baku banyaknya pelanggan yang
dilayani kedai kopi ini dalam satu jam pada waktu seperti ini?
\item Apakah kedatangan hanya 60 pelanggan ke kedai kopi ini dalam
satu jam pada waktu seperti ini dapat dianggap sangat rendah hingga tidak biasa?
\item Hitunglah peluang bahwa kedai kopi ini melayani 70 pelanggan dalam
satu jam pada waktu seperti ini.
\end{parts}
}{}

% 32

\eoce{\qt{Salah ketik seorang stenografer\label{stenographer_typos}} Seorang stenografer
pengadilan yang sangat terampil rata-rata membuat satu salah ketik per jam.
\begin{parts}
\item Distribusi peluang manakah yang paling tepat untuk menghitung
peluang bahwa stenografer ini membuat sejumlah tertentu salah ketik dalam satu jam?
\item Berapakah rata-rata dan simpangan baku banyaknya salah ketik
yang dibuat stenografer ini dalam satu jam?
\item Apakah akan dianggap tidak biasa jika stenografer ini membuat 4 salah ketik dalam
satu jam tertentu?
\item Hitunglah peluang bahwa stenografer ini membuat paling banyak 2 salah ketik
dalam satu jam tertentu.
\end{parts}
}{}

% 33

\eoce{\qtq{Berapa banyak mobil yang datang\label{cars_in_parking_lot}}
Pada hari Senin hingga Kamis yang bukan hari libur,
rata-rata banyaknya kendaraan yang mengunjungi sebuah
toko antara pukul 14.00 dan 15.00 setiap sore adalah 6.5,
dan banyaknya mobil yang datang pada suatu hari
mengikuti distribusi Poisson.
\begin{parts}
\item
    Berapakah peluang bahwa tepat
    5 mobil akan datang pada Senin depan?
\item
    Berapakah peluang bahwa
    0, 1, atau 2 mobil akan datang pada Senin depan
    antara pukul 14.00 dan 15.00?
\item
    Rata-rata ada 11.7 orang yang datang
    dengan kendaraan pada jam yang sama.
    Apakah banyaknya orang yang datang
    dengan mobil selama satu jam tersebut juga kemungkinan mengikuti distribusi Poisson?
    Jelaskan.
\end{parts}
}{}

% 34

\eoce{\qt{Bagasi hilang\label{lost_baggage}}
Sesekali, sebuah maskapai kehilangan bagasi.
Misalkan sebuah maskapai kecil mendapati bahwa banyaknya
bagasi yang hilang setiap hari kerja dapat dimodelkan dengan cukup baik menggunakan
model Poisson dengan rata-rata 2.2 bagasi.
\begin{parts}
\item
    Berapakah peluang bahwa maskapai tersebut
    tidak akan kehilangan bagasi pada Senin depan?
\item
    Berapakah peluang bahwa maskapai tersebut
    akan kehilangan 0, 1, atau 2 bagasi pada Senin depan?
\item
    Misalkan maskapai tersebut berkembang selama
    3 tahun berikutnya sehingga jumlah penerbangannya menjadi dua kali lipat,
    dan direktur utamanya bertanya apakah
    mereka masih wajar melanjutkan penggunaan
    model Poisson dengan rata-rata~2.2.
    Rekomendasi apakah yang tepat?
    Jelaskan.
\end{parts}
}{}
